\chapter{PMS and PMDD}

\section*{Definition}
Premenstrual syndrome (PMS) is a recurrent luteal-phase pattern of emotional, behavioral, and physical symptoms that resolve soon after menses onset. Premenstrual dysphoric disorder (PMDD) is a more severe form characterized by marked mood symptoms (irritability, depression, anxiety, affective lability) that cause clinically significant impairment and meet DSM-5 criteria.

\section*{What Happens in Your Body}
PMS/PMDD involves an abnormal central nervous system response to normal cyclic hormonal fluctuations. Altered serotonergic and GABAergic signaling in the luteal phase (from progesterone metabolite allopregnanolone) causes emotional dysregulation. Progesterone withdrawal triggers neurosteroid-mediated GABAA receptor changes, reducing inhibitory tone and increasing anxiety and irritability.

\section*{Causes}
\begin{itemize}
  \item \textbf{Hormonal:} Exaggerated sensitivity to normal progesterone and estrogen fluctuations; altered allopregnanolone modulation of GABAA receptors
  \item \textbf{Genetic:} Heritability ~50\%; polymorphisms in serotonin transporter (5-HTTLPR) and estrogen receptor genes
  \item \textbf{Neurotransmitter:} Reduced luteal-phase serotonergic transmission, altered beta-endorphin levels
  \item \textbf{Risk factors:} History of major depressive disorder, perinatal depression, high stress, interpersonal trauma, family history of PMS/PMDD
  \item \textbf{Nutritional:} Low calcium, magnesium, and vitamin D intake may contribute
\end{itemize}

\section*{When to See a Doctor}
See your doctor if:
\begin{itemize}
  \item Symptoms consistently occur in the week before menses and remit after onset, causing significant interpersonal or occupational impairment
  \item Mood symptoms include suicidal ideation, severe depression, or anger that feels uncontrollable
  \item Symptoms do not resolve with lifestyle changes and over-the-counter measures
\end{itemize}

\section*{Related Symptoms}
\begin{itemize}
  \item Irritability, mood swings, depression, anxiety, bloating, breast tenderness, fatigue, sleep disturbance, food cravings, headache, joint pains, difficulty concentrating
\end{itemize}
\textbf{Important:} A prospective symptom diary charted daily across two cycles is essential to distinguish PMS/PMDD from other mood disorders.

\section*{Self-Care / Home Management}
\begin{itemize}
  \item Exercise regularly (aerobic exercise reduces symptoms)
  \item Dietary changes: reduce sodium, caffeine, and refined sugar; increase complex carbohydrates, calcium (1000--1200 mg/day), and vitamin B6
  \item Stress management: yoga, mindfulness, relaxation techniques
  \item Over-the-counter NSAIDs (ibuprofen, naproxen) for pain and bloating
\end{itemize}

\section*{How Your Doctor Will Evaluate You}
\begin{itemize}
  \item Prospective symptom charting (daily) for 2 full cycles -- the gold standard
  \item For PMDD: DSM-5 criteria requiring at least 5 symptoms (including 1 mood symptom) in the luteal phase with subsidence after menses
  \item PHQ-9 and GAD-7 to screen for comorbid depression and anxiety
  \item Rule out thyroid dysfunction, anemia, and perimenopause
\end{itemize}

\section*{Treatment Options}
\begin{itemize}
  \item First-line pharmacotherapy for PMDD: SSRIs (fluoxetine, sertraline, citalopram) -- can be dosed continuously or only in the luteal phase
  \item Combined oral contraceptives: drospirenone/ethinyl estradiol (Yaz, Yasmin) shown effective for PMDD
  \item GnRH agonists (leuprolide) with add-back hormone therapy for severe, refractory cases
  \item CBT for PMS/PMDD -- effective for coping with emotional symptoms
  \item Severe: gonadotropin-releasing hormone antagonist or ovariectomy (rarely needed)
\end{itemize}

\section*{References}
\begin{thebibliography}{99}
\bibitem{pms_pmdd:ref1} Yonkers KA, Simoni MK. ``Premenstrual disorders.'' \textit{American Journal of Obstetrics and Gynecology}, 2018;218(1):68--74.
\bibitem{pms_pmdd:ref2} Ismaili E, Walsh S, O'Brien PMS, et al. ``The management of premenstrual syndrome.'' \textit{Clinical Obstetrics and Gynecology}, 2019;62(2):327--341.
\bibitem{pms_pmdd:ref3} Pearlstein T, Steiner M. ``Premenstrual dysphoric disorder: burden of illness and treatment update.'' \textit{Journal of Psychiatry \& Neuroscience}, 2008;33(4):291--301.
\bibitem{pms_pmdd:ref4} Land\'en M, Eriksson E. ``How does premenstrual dysphoric disorder relate to depression and anxiety disorders?'' \textit{Depression and Anxiety}, 2003;17(3):122--129.
\bibitem{pms_pmdd:ref5} Shah NR, Jones JB, Aperi J, et al. ``Selective serotonin reuptake inhibitors for premenstrual syndrome and premenstrual dysphoric disorder.'' \textit{Annals of Pharmacotherapy}, 2008;42(7):1036--1042.
\bibitem{pms_pmdd:ref6} Freeman EW. ``Premenstrual syndrome and premenstrual dysphoric disorder: definitions and diagnosis.'' \textit{Psychoneuroendocrinology}, 2003;28(Suppl 3):25--37.

\end{thebibliography}
